Considera la funció
\[
f(x)=\frac{x+3}{x-1} .
\]

\begin{apartats}

\apartat[1,25]{1}
Troba l'equació de la recta tangent a la gràfica de $f$ en el punt d'abscissa $x=3$.

\begin{solucio}
$f(3)=\dfrac{6}{2}=3$. $f'(x)=\dfrac{(x-1)-(x+3)}{(x-1)^2}=\dfrac{-4}{(x-1)^2}$, d'on
$f'(3)=-1$.\\
Tangent: $y-3=-(x-3)$, és a dir $y=-x+6$.
\end{solucio}

\apartat[1,25]{0,75}
Calcula l'àrea del triangle que aquesta recta tangent forma amb els eixos de coordenades.

\begin{solucio}
La recta $y=-x+6$ talla els eixos en $(6,0)$ i $(0,6)$. El triangle és rectangle, amb catets
de longitud $6$: l'àrea és $\dfrac{6\cdot6}{2}=\boxed{18}$ unitats quadrades.
\end{solucio}

\begin{nomesllarg}
\apartat{0,75}
Hi ha algun altre punt de la gràfica de $f$ en què la recta tangent sigui paral·lela a
aquesta? Si n'hi ha, troba'l.

\begin{solucio}
Cal $f'(x)=-1$: $\dfrac{-4}{(x-1)^2}=-1\iff(x-1)^2=4\iff x=3$ o $x=-1$.\\
Sí: el punt d'abscissa $x=-1$, que és $\left(-1,f(-1)\right)=(-1,-1)$.
\end{solucio}
\end{nomesllarg}

\end{apartats}
